\clearpage
\begin{explanationbox}
\textbf{几何直觉一：差固定，按比例分份，求出两边之和。}

把焦点三角形 $\triangle PF_1F_2$ 的两条腰记作
$r_1=PF_1$、$r_2=PF_2$。双曲线定义规定
\[
\boxed{|r_1-r_2|=2a},\qquad \boxed{F_1F_2=2c}.
\]
如果题目给 $r_1:r_2=\lambda:1$，就把两条腰分别看成
“$\lambda$ 份”和“$1$ 份”：\textbf{两腰的差 $|\lambda-1|$ 份恰好等于 $2a$}。
于是
\[
\text{一份}=\frac{2a}{|\lambda-1|},\qquad
r_1+r_2=\frac{2a(\lambda+1)}{|\lambda-1|}.
\]
但三角形两边之和不能小于第三边 $2c$：
\[
\frac{2a(\lambda+1)}{|\lambda-1|}\ge2c
\ \Longrightarrow\ \boxed{e\le\frac{\lambda+1}{|\lambda-1|}}.
\]
\textbf{理解优先级：}不必先背最后的分式。只记“\textbf{差 $2a$、底边 $2c$、比值分份}”，
公式即可现推。
\end{explanationbox}

\clearpage
\begin{explanationbox}
\textbf{几何直觉二：为什么最值恰在顶点？}

当 $P$ 是右顶点 $A=(a,0)$，三个点共线，且 $P$ 位于两焦点之间。
三角形退化为一条线段，所以
\[
\boxed{PF_1+PF_2=F_1F_2=2c}.
\]
此时刚好达到不等式的\textbf{等号}；离开顶点，真正组成三角形，便有
$PF_1+PF_2>2c$，所以同一个 $\lambda$ 下 $e$ 必须\textbf{严格小于}上界。

\begin{center}\begin{tikzpicture}[font=\small,>=stealth]
\draw[gray!45,rounded corners=3pt] (0,0) rectangle (14.2,8.5);
\node at (7.1,7.85) {等号就在顶点：三点共线，两腰恰好拼成焦距};

\begin{scope}[shift={(6,5.8)},scale=1.05]
\draw[gray] (-3.5,0)--(3.5,0);
\draw[orange!85!black,very thick] (-3,0)--(1,0);
\draw[purple!80!black,very thick] (1,0)--(3,0);
\fill (-3,0) circle(2.8pt) node[below=6pt] {$F_1$};
\fill (3,0) circle(2.8pt) node[below=6pt] {$F_2$};
\fill (1,0) circle(3pt) node[below=6pt] {$A_{+}$};
\node[orange!85!black,above=6pt] at (-1.0,0) {$r_1=4$};
\node[purple!80!black,above=6pt] at (2.0,0) {$r_2=2$};
\end{scope}
\node[anchor=west,align=left] at (10.4,5.8) {右顶点\\[10pt]$\lambda=2$};

\begin{scope}[shift={(6,3.0)},scale=1.05]
\draw[gray] (-3.5,0)--(3.5,0);
\draw[orange!85!black,very thick] (-3,0)--(-1,0);
\draw[purple!80!black,very thick] (-1,0)--(3,0);
\fill (-3,0) circle(2.8pt) node[below=6pt] {$F_1$};
\fill (3,0) circle(2.8pt) node[below=6pt] {$F_2$};
\fill (-1,0) circle(3pt) node[below=6pt] {$A_{-}$};
\node[orange!85!black,above=6pt] at (-2.0,0) {$r_1=2$};
\node[purple!80!black,above=6pt] at (1.0,0) {$r_2=4$};
\end{scope}
\node[anchor=west,align=left] at (10.4,3.0) {左顶点\\[10pt]$\lambda=1/2$};
\node at (7.1,.65) {两行均取 $a=1,c=3,e=3$；两腰之和均为 $6=2c$。};
\end{tikzpicture}
\end{center}

\textbf{几何直觉三：阿波罗尼斯圆（定比轨迹），不是“两焦点圆”。}
固定 $\lambda>0$，把\emph{平面上所有}满足 $PF_1=\lambda PF_2$ 的点收集起来：
$\lambda\ne1$ 时构成\textbf{一条圆} $\mathcal A_\lambda$。其圆心在焦点连线所在的 $x$ 轴上：
\[
\boxed{C_\lambda\!\left(\frac{c(1+\lambda^2)}{\lambda^2-1},0\right)},\qquad
\boxed{R_\lambda=\frac{2c\lambda}{|\lambda^2-1|}}.
\]
\textbf{推导：}把 $PF_1^2=\lambda^2PF_2^2$ 展开、配方：
\[
(x+c)^2+y^2=\lambda^2[(x-c)^2+y^2]
\iff
(x-x_C)^2+y^2=R_\lambda^2.
\]
因此 C086 的问题实质上是：\textbf{圆 $\mathcal A_\lambda$ 能否与双曲线相交？}
例如 $\lambda=2$、$a=2,c=4$ 时，圆与右支相交于上下对称的两个点。

\begin{center}% C086 阿波罗尼斯圆与双曲线精确示意（数学坐标，非自由重绘）
% H:x^2/4-y^2/12=1, a=2, c=4, e=2; lambda=2.
% F1=(-4,0), F2=(4,0).
% Apollonius circle center C=(20/3,0), radius R=16/3.
% H cap A=(2,0); circle leftmost T=(4/3,0).
% Exact intersections P+/-=(3,+/-sqrt(15)); PF1=8, PF2=4.
\begin{tikzpicture}[x=0.58cm,y=0.58cm,>=stealth]
  \draw[->,gray!65] (-4.85,0)--(12.55,0) node[right] {$x$};
  \draw[->,gray!65] (0,-5.8)--(0,5.8) node[above] {$y$};
  % Circle locus. The path is constructed from exact center and radius.
  \draw[orange!85!black,very thick] ({20/3},0) circle [radius={16/3}];
  % Entire visible hyperbola branches from analytic formula y=+/-sqrt(3(x^2-4)).
  \draw[teal!80!black,very thick,samples=95,domain=2:3.75,smooth]
     plot (\x,{sqrt(3*(\x*\x-4))});
  \draw[teal!80!black,very thick,samples=95,domain=2:3.75,smooth]
     plot (\x,{-sqrt(3*(\x*\x-4))});
  \draw[teal!80!black,very thick,samples=95,domain=-3.75:-2,smooth]
     plot (\x,{sqrt(3*(\x*\x-4))});
  \draw[teal!80!black,very thick,samples=95,domain=-3.75:-2,smooth]
     plot (\x,{-sqrt(3*(\x*\x-4))});
  \coordinate (Fone) at (-4,0);
  \coordinate (Ftwo) at (4,0);
  \coordinate (Ap) at (2,0);
  \coordinate (Cc) at ({20/3},0);
  \coordinate (T) at ({4/3},0);
  \coordinate (Pup) at (3,{sqrt(15)});
  \coordinate (Pdn) at (3,{-sqrt(15)});
  \draw[dashed,purple!75!black,thin] (Fone)--(Pup)--(Ftwo);
  \fill[black] (Fone) circle (2pt) node[below left] {$F_1$};
  \fill[black] (Ftwo) circle (2pt) node[below=5pt] {$F_2$};
  \fill[teal!75!black] (Ap) circle (2.3pt) node[below right=10pt] {$A$};
  \fill[orange!85!black] (Cc) circle (2pt) node[below=5pt] {$C_\lambda$};
  \fill[orange!85!black] (T) circle (2.2pt) node[above left=8pt] {$T$};
  \fill[red!80!black] (Pup) circle (2.8pt) node[above left] {$P_+$};
  \fill[red!80!black] (Pdn) circle (2.8pt) node[below left] {$P_-$};
  \node[fill=white,fill opacity=0.9,text opacity=1,text=orange!80!black]
    at (8.7,5.65) {$\mathcal A_2$: 定比轨迹圆};
  \node[fill=white,fill opacity=0.9,text opacity=1,text=teal!75!black]
    at (-1.3,5.35) {$H$: 定差轨迹};
  \node[fill=white,fill opacity=0.9,text opacity=1] at (8.6,-4.75)
    {$P_\pm=(3,\pm\sqrt{15})$};
\end{tikzpicture}
\end{center}

\textbf{最直观的圆上长度变化（先看 $\lambda>1$）：}
右焦点 $F_2$ 在圆内；当 $P$ 沿 $\mathcal A_\lambda$ 走一圈，
$PF_2$ 最小值和最大值分别是
\[
\boxed{\min PF_2=\frac{2c}{\lambda+1}},\qquad
\boxed{\max PF_2=\frac{2c}{\lambda-1}}.
\]
因 $PF_1-PF_2=(\lambda-1)PF_2$，所以距离差在
\[
\boxed{\frac{2c(\lambda-1)}{\lambda+1}}
\ \le\ PF_1-PF_2\ \le\ 2c
\]
之间连续变化。双曲线要求\textbf{距离差恰好 $2a$}；由于 $2a<2c$，
存在交点当且仅当
$2a\ge 2c(\lambda-1)/(\lambda+1)$，即
\[
\boxed{e=\frac ca\le\frac{\lambda+1}{\lambda-1}}.
\]
\textbf{等号的几何意义：}圆上距离差取最小值的点在 $x$ 轴上，
且此时正好是双曲线的右顶点；严格不等时有上、下两个交点。
当 $0<\lambda<1$，交换 $F_1,F_2$（或将 $\lambda$ 换成 $1/\lambda$）得到同样的绝对值形式。

\textbf{不要混淆：}这里的\textbf{一条阿波罗尼斯圆}是全部定比点的轨迹；
证明存在性时还可以另作\textbf{两个以焦点为圆心的辅助圆}（见证明方法五）。
\end{explanationbox}
